
\SourceLeaf{476}

\Needspace{5\baselineskip}\begin{center}INHALT.\end{center}
\BookRule{0.15\linewidth}

\ContentsHeading{\textsc{Widmung}}{S. 3.}

\ContentsHeading{\textsc{Vorrede}}{S. 5.}

\ContentsHeading{\textsc{Erster Abschnitt.} Die Kongruenz von Zahlen im Allgemeinen}{S. 9.}

\ContentsEntry{Kongruente Zahlen, Moduln, Reste und Nichtreste, Artikel 1 f.}

\ContentsEntry{Kleinste Reste, 4.}

\ContentsEntry{Elementare Sätze über kongruente Zahlen, 5.}

\ContentsEntry{Einige Anwendungen, 12.}

\ContentsHeading{\textsc{Zweiter Abschnitt.} Kongruenzen ersten Grades}{S. 14.}

\ContentsEntry{Vorbereitende Sätze über Primzahlen, Faktoren usw., Artikel 13.}

\ContentsEntry{Lösung der Kongruenzen ersten Grades, 26.}

\ContentsEntry{Bestimmung einer Zahl, die nach gegebenen Moduln gegebenen Resten kongruent ist, 32.}

\ContentsEntry{Lineare Kongruenzen mit mehreren Unbekannten, 37.}

\ContentsEntry{Verschiedene Sätze, 38.}

\ContentsHeading{\textsc{Dritter Abschnitt.} Reste von Potenzen}{S. 38.}

\ContentsEntry{Die Reste der Glieder einer mit der Einheit beginnenden geometrischen Folge bilden eine periodische Folge, Artikel 45.}

\ContentsEntry{\textit{Zuerst werden Primzahlmoduln betrachtet.}}

\ContentsEntry{Setzt man den Modul \(=p\), so teilt die Anzahl der Glieder der Periode die Zahl \(p-1\), Artikel 49.}

\ContentsEntry{Der Fermatsche Satz, 50.}

\ContentsEntry{Wie vielen Zahlen Perioden entsprechen, deren Gliederzahl ein gegebener Teiler von \(p-1\) ist, Artikel 52.}

\ContentsEntry{Primitive Wurzeln, Basen, Indizes, 57.}

\ContentsEntry{Der Algorithmus der Indizes, 58.}

\ContentsEntry{Die Wurzeln der Kongruenz \(x^n\equiv A\), Artikel 60.}

\ContentsEntry{Der Zusammenhang der Indizes in verschiedenen Systemen, 69.}

\ContentsEntry{Für besondere Zwecke geeignete Basen, 72.}

\ContentsEntry{Eine Methode zur Bestimmung primitiver Wurzeln, 73.}

\ContentsEntry{Verschiedene Sätze über Perioden und primitive Wurzeln, 75.}

\ContentsEntry{(Der Wilsonsche Satz, 76).}

\SourceLeaf{477}

\ContentsEntry{\textit{Moduln, die Potenzen von Primzahlen sind}, Artikel 82.}

\ContentsEntry{\textit{Moduln, die Zweierpotenzen sind}, 90.}

\ContentsEntry{\textit{Aus mehreren Primzahlen zusammengesetzte Moduln}, 92.}

\ContentsHeading{\textsc{Vierter Abschnitt.} Kongruenzen zweiten Grades}{S. 73.}

\ContentsEntry{Quadratische Reste und Nichtreste, Artikel 94.}

\ContentsEntry{Ist der Modul eine Primzahl, so ist die Anzahl der Reste, die kleiner als er sind, gleich der Anzahl der Nichtreste, 96.}

\ContentsEntry{Ob eine zusammengesetzte Zahl Rest oder Nichtrest einer gegebenen Primzahl ist, hängt von der Beschaffenheit ihrer Faktoren ab, 98.}

\ContentsEntry{Zusammengesetzte Moduln, 100.}

\ContentsEntry{Ein allgemeines Kriterium dafür, ob eine gegebene Zahl Rest oder Nichtrest einer gegebenen Primzahl ist, 106.}

\ContentsEntry{\textit{Untersuchungen über Primzahlen, für die gegebene Zahlen Reste oder Nichtreste sind}, 107.}

\ContentsEntry{Der Rest \(-1\), Artikel 108.}

\ContentsEntry{Die Reste \(+2\) und \(-2\), Artikel 112.}

\ContentsEntry{Die Reste \(+3\) und \(-3\), Artikel 117.}

\ContentsEntry{Die Reste \(+5\) und \(-5\), Artikel 121.}

\ContentsEntry{Über \(\pm7\), Artikel 124.}

\ContentsEntry{Vorbereitung der allgemeinen Untersuchung, 125.}

\ContentsEntry{Durch Induktion wird ein allgemeiner (\textit{fundamentaler}) Satz aufgestellt, und Folgerungen werden daraus gezogen, 130.}

\ContentsEntry{Ein strenger Beweis dieses Satzes, 135.}

\ContentsEntry{Eine entsprechende Methode zum Beweis des Satzes aus Artikel 114, 145.}

\ContentsEntry{Lösung des allgemeinen Problems, 146.}

\ContentsEntry{Lineare Formen, die alle Primzahlen enthalten, für die eine beliebige gegebene Zahl Rest oder Nichtrest ist, 147.}

\ContentsEntry{Die Arbeiten anderer über diese Untersuchungen, 151.}

\ContentsEntry{Nicht reine Kongruenzen zweiten Grades, 152.}

\ContentsHeading{\textsc{Fünfter Abschnitt.} Formen und unbestimmte Gleichungen zweiten Grades}{S. 120.}

\ContentsEntry{Das Ziel der Untersuchung; Definition und Bezeichnung der Formen, Artikel 153.}

\ContentsEntry{Darstellung von Zahlen; die Determinante, 154.}

\ContentsEntry{Werte des Ausdrucks \(\surd(bb-ac)\) (mod. \(M\)), zu denen die Darstellung der Zahl \(M\) durch die Form \((a,b,c)\) gehört, 155.}

\ContentsEntry{Eine Form, die eine andere enthält oder unter einer anderen enthalten ist; eigentliche und uneigentliche Transformationen, 157.}

\ContentsEntry{Eigentliche und uneigentliche Äquivalenz, 158.}

\ContentsEntry{Entgegengesetzte Formen, 159, benachbarte Formen, 160.}

\ContentsEntry{Gemeinsame Teiler der Koeffizienten von Formen, 161.}

\ContentsEntry{Der Zusammenhang aller gleichartigen Transformationen einer gegebenen Form in eine gegebene Form, 162.}

\ContentsEntry{Ambige Formen, 163.}

\ContentsEntry{Ein Satz über den Fall, dass eine Form unter einer anderen zugleich eigentlich und uneigentlich enthalten ist, 164.}

\ContentsEntry{Allgemeines über Darstellungen von Zahlen durch Formen und ihren Zusammenhang mit Transformationen, 166.}

\ContentsEntry{\textit{Formen negativer Determinante}, 171.}

\ContentsEntry{Besondere Anwendungen auf die Zerlegung von Zahlen in zwei Quadrate, in ein Quadrat und das Doppelte eines Quadrats sowie in ein Quadrat und das Dreifache eines Quadrats, 182.}

\ContentsEntry{\textit{Formen positiver, nichtquadratischer Determinante}, 183.}

\SourceLeaf{478}

\ContentsEntry{\textit{Formen mit quadratischer Determinante}, Artikel 206.}

\ContentsEntry{Formen, die unter anderen enthalten, ihnen aber nicht äquivalent sind, 213.}

\ContentsEntry{\textit{Formen der Determinante \(0\)}, Artikel 215.}

\ContentsEntry{Die allgemeine Lösung sämtlicher unbestimmter Gleichungen zweiten Grades mit zwei Unbekannten durch ganze Zahlen, 216.}

\ContentsEntry{Historische Anmerkungen, 222.}

\Needspace{5\baselineskip}\begin{center}\textit{Weitere Untersuchungen über Formen.}\end{center}

\ContentsEntry{Einteilung der Formen gegebener Determinante in Klassen, Artikel 223; der Klassen in Ordnungen, 226.}

\ContentsEntry{Einteilung der Ordnungen in Geschlechter, 228.}

\ContentsEntry{\textit{Komposition der Formen}, 234.}

\ContentsEntry{Komposition der Ordnungen, 245, der Geschlechter, 246, der Klassen, 249.}

\ContentsEntry{Bei gegebener Determinante enthalten die einzelnen Geschlechter derselben Ordnung gleich viele Klassen, 252.}

\ContentsEntry{Vergleich der Anzahlen von Klassen in den einzelnen Geschlechtern verschiedener Ordnungen, 253.}

\ContentsEntry{Die Anzahl der ambigen Klassen, 257.}

\ContentsEntry{Mindestens der Hälfte aller für eine gegebene Determinante angebbaren Charaktere können keine eigentlich primitiven Geschlechter entsprechen (bei negativer Determinante positive), 261.}

\ContentsEntry{Ein zweiter Beweis des Fundamentaltheorems und der übrigen Sätze über die Reste \(-1,\allowbreak +2,\allowbreak -2\), 262.}

\ContentsEntry{Die Hälfte der Charaktere, denen keine Geschlechter entsprechen können, wird näher bestimmt, 263.}

\ContentsEntry{Eine besondere Methode zur Zerlegung von Primzahlen in zwei Quadrate, 265.}

\Needspace{5\baselineskip}\begin{center}\textit{Ein Exkurs mit einer Abhandlung über ternäre Formen}, Artikel 266 ff.\end{center}

\ContentsEntry{\textit{Einige Anwendungen auf die Theorie der binären Formen.}}

\ContentsEntry{Bestimmung einer Form, aus deren Verdopplung eine gegebene binäre Form des Hauptgeschlechts entsteht, 286.}

\ContentsEntry{Allen Charakteren außer den in den Artikeln 262, 263 als unmöglich erkannten entsprechen tatsächlich Geschlechter, 287 III.}

\ContentsEntry{Die Theorie der Zerlegung sowohl von Zahlen als auch von binären Formen in drei Quadrate, 288.}

\ContentsEntry{Beweis der Fermatschen Sätze, dass jede ganze Zahl in drei Dreieckszahlen oder vier Quadrate zerlegt werden kann, 293.}

\ContentsEntry{Lösung der Gleichung \(axx+byy+czz=0\), Artikel 294.}

\ContentsEntry{Die Methode, mit der der berühmte Le Gendre das Fundamentaltheorem behandelte, 296.}

\ContentsEntry{Darstellung der Null durch beliebige ternäre Formen, 299.}

\ContentsEntry{Die allgemeine Lösung unbestimmter Gleichungen zweiten Grades mit zwei Unbekannten durch rationale Größen, 300.}

\ContentsEntry{Die durchschnittliche Anzahl der Geschlechter. 301; der Klassen, 302.}

\ContentsEntry{Ein besonderer Algorithmus der eigentlich primitiven Klassen; reguläre und irreguläre Determinanten usw., 305.}

\ContentsHeading{\textsc{Sechster Abschnitt.} Verschiedene Anwendungen der vorangegangenen Untersuchungen}{S. 380.}

\ContentsEntry{Zerlegung von Brüchen in einfachere, Artikel 309.}

\ContentsEntry{Umwandlung gewöhnlicher Brüche in Dezimalbrüche, 312.}

\ContentsEntry{Lösung der Kongruenz \(xx\equiv A\) durch die Ausschlussmethode, 319.}

\ContentsEntry{Lösung der unbestimmten Gleichung \(mxx+nyy=A\) durch Ausschlüsse, 323.}

\SourceLeaf{479}

\ContentsEntry{Eine andere Methode zur Lösung der Kongruenz \(xx\equiv A\) für den Fall, dass \(A\) negativ ist, Artikel 327.}

\ContentsEntry{Zwei Methoden, zusammengesetzte Zahlen von Primzahlen zu unterscheiden und die Faktoren der Ersteren zu bestimmen, 329.}

\ContentsHeading{\textsc{Siebter Abschnitt.} Gleichungen, welche die Kreisteilung bestimmen}{S. 412.}

\ContentsEntry{Die Untersuchung wird auf den einfachsten Fall zurückgeführt, in dem die Anzahl der Teile, in die der Kreis geteilt werden soll, eine Primzahl ist, Artikel 336.}

\ContentsEntry{Gleichungen für trigonometrische Funktionen von Bögen, die ein Teil oder mehrere Teile des ganzen Umfangs sind; Zurückführung der trigonometrischen Funktionen auf Wurzeln der Gleichung \(x^n-1=0\), Artikel 337.}

\ContentsEntry{\textit{Theorie der Wurzeln dieser Gleichung} (wobei vorausgesetzt wird, dass \(n\) eine Primzahl ist).}

\ContentsEntry{Lässt man die Wurzel \(1\) weg, so sind die übrigen Wurzeln \((\Omega)\) in der Gleichung \(X=x^{n-1}+x^{n-2}+\mathrm{usw.}+x+1=0\) enthalten.}

\ContentsEntry{Die Funktion \(X\) lässt sich nicht in Faktoren niedrigeren Grades zerlegen, deren Koeffizienten sämtlich rational sind, 341.}

\ContentsEntry{Das Ziel der folgenden Untersuchungen wird dargelegt, 342.}

\ContentsEntry{Alle Wurzeln \(\Omega\) werden in bestimmte Klassen (Perioden) eingeteilt, 343.}

\ContentsEntry{Verschiedene Sätze über diese Perioden, 344 ff.}

\ContentsEntry{Auf diese Untersuchungen wird die Lösung der Gleichung \(X=0\) gegründet, Artikel 352.}

\ContentsEntry{Beispiele für \(n=19\), wobei die Aufgabe auf zwei kubische Gleichungen und eine quadratische zurückgeführt wird, und für \(n=17\), wobei sie auf vier quadratische zurückgeführt wird, Artikel 353. 354.}

\Needspace{4\baselineskip}\ContentsEntry{\textit{Weitere Untersuchungen über diesen Gegenstand.}}

\ContentsEntry{Summen mit einer geraden Anzahl von Gliedern sind reelle Größen, 355.}

\ContentsEntry{Die Gleichung, welche die Einteilung der Wurzeln \(\Omega\) in zwei Perioden bestimmt, 356.}

\ContentsEntry{Beweis des im vierten Abschnitt erwähnten Satzes, 357.}

\ContentsEntry{Die Gleichung zur Einteilung der Wurzeln \(\Omega\) in drei Perioden, 358.}

\ContentsEntry{Zurückführung der Gleichungen, durch die die Wurzeln \(\Omega\) gefunden werden, auf reine Gleichungen, 359.}

\ContentsEntry{\textit{Anwendung der vorangegangenen Untersuchungen auf trigonometrische Funktionen.}}

\ContentsEntry{Eine Methode zur Bestimmung der Winkel, denen die einzelnen Wurzeln \(\Omega\) entsprechen, 361.}

\ContentsEntry{Tangens, Kotangens, Sekans und Kosekans werden ohne Division aus Sinus und Kosinus abgeleitet, 362.}

\ContentsEntry{Eine Methode zur schrittweisen Senkung des Grades der Gleichungen für trigonometrische Funktionen, 363.}

\ContentsEntry{Kreisteilungen, die sich durch quadratische Gleichungen oder durch geometrische Konstruktionen ausführen lassen, 365.}

\ContentsHeading{\textsc{Nachträge}}{S. 465.}

\ContentsHeading{\textsc{Tabellen}}{S. 467.}

\BookRule{0.3\linewidth}

\Needspace{5\baselineskip}\begin{center}\textsc{Gedruckt in Goslar bei E. W. G. Kircher.}\end{center}
\clearpage
